% ======================================================================
%  Section 4.1 --- Introduction.  Corrected.
%
%  SIX CHANGES, in order of importance. Each is explained at the point it
%  occurs in the text below, in a %%% comment.
%
%  1. The opening sentence claimed Chapter 3 shows convergence "accurately"
%     under partial occlusion. Section 4.7.2 shows that same scheme
%     converging 26 pixels from the desired pose. The chapter's own results
%     contradict its first sentence.
%  2. "The Hermite bank ALREADY COMPUTED by the scheme carries exactly that
%     missing information" is no longer true. After the correction to 4.4.2
%     the measure is built from the first-order kernels D_{1,0}, D_{0,1},
%     which are NOT among the summed kernels the feature vector uses; the
%     implementation adds six separable convolutions per iteration.
%  3. "A three-way ablation ANSWERS IT DIRECTLY" promises a verdict that
%     4.7.2 does not deliver: Variants 2 and 3 differ by under eight per
%     cent, in opposite directions.
%  4. \label{sec:ch5-intro} in a Chapter 4 file.
%  5. The section list switches from semicolons to commas half way through,
%     and points at sec:ablation-results for "the results" when the results
%     occupy the whole of 4.7.
%  6. Spelling: organized/characterizes are -ize, while the chapter body
%     uses -ise elsewhere. Pick one for the whole thesis.
% ======================================================================

\section{Introduction}
\label{sec:ch4-intro}
%%% (4) was sec:ch5-intro. Any \ref to the old label elsewhere must follow.

%%% (1) "converges accurately under partial occlusion" is too strong, and
%%% Section 4.7.2 disproves it: the same scheme, Variant 1, terminates
%%% 1.29 mm and 1.70 degrees from the desired pose, a misalignment of 26
%%% pixels, with its commanded velocity already vanished. What Chapter 3
%%% established is that the scheme remains stable and converges; what it did
%%% not establish is that it converges to the right pose. Saying so here is
%%% what motivates the chapter, rather than undercutting it.
Chapter~\ref{ch:herpvs} showed that HER-PVS retains its convergence under
partial occlusion, and Section~\ref{sec:hermite-properties} established the
structural reason: the Hermite filters have compact support, so an occlusion
corrupts only those responses whose support overlaps the occluded region.

That tolerance is nevertheless \emph{passive}. The corrupted responses are
diluted within a dense feature vector but never identified, and they continue
to enter the product $\mathbf{L}_{R}^{\top}\mathbf{e}$ that determines the
commanded velocity. Because the occluded region is a coherent block of large,
correlated residuals rather than a random perturbation, it biases the descent
direction; and because dilution is an argument about proportions, its
protection degrades as the occluded fraction grows. Passive tolerance
therefore offers no graceful degradation, and, as
Section~\ref{sec:ablation-results} shows, it does not prevent the scheme from
converging to a systematically displaced pose.
%%% (1, continued) This forward reference replaces the accuracy claim that
%%% was removed from the first sentence. The bias is no longer asserted in
%%% the abstract; it is measured in 4.7.2, and the introduction says where.

This chapter makes the treatment of occlusion \emph{active}, attenuating the
influence of corrupted measurements through a weighting matrix acting on both
the error and the interaction matrix, within the framework of robust
M-estimation~\cite{key26}. The contribution lies not in the use of an
M-estimator as such, but in the information from which the weights are
derived. A standard M-estimator is \emph{structure-blind}: it assesses each
measurement by the magnitude of its own residual alone, and knows nothing of
what the image contains at that position. The Hermite representation on which
the scheme already operates supplies exactly that missing information, since
its local response characterises the image structure at the very scales on
which the control acts. The measure used here is drawn from that
representation --- the first-order kernels of the same Gauss--Hermite family,
evaluated at the same three analysis scales as the features themselves --- at
the cost of six additional separable convolutions per iteration. The resulting
weighting is accordingly termed \emph{Hermite-informed}.
%%% (2) The previous wording said the bank "already computed by the scheme"
%%% carries the information, implying the measure is free. It is not. The
%%% feature vector uses the SUMMED kernels; the structure measure of (4.10)
%%% uses the first-order kernels, which the implementation evaluates
%%% separately. The claim that survives is the one made above: same family,
%%% same scales, six extra convolutions. The same correction is needed in
%%% 4.4.1, p. 67, which says "the only additional work being a pointwise
%%% energy evaluation and a mean over the image".

%%% (3) "answers it directly" was a promise the results do not keep: under
%%% occlusion Variants 2 and 3 differ by 7.6 per cent in rotation and 7.2
%%% per cent in translation, in opposite directions. The ablation is still
%%% the right instrument and is still worth describing here; what it must
%%% not do in the introduction is announce a verdict.
Whether this structural information is worth its complexity is an empirical
question, and the chapter is built around an experiment designed to isolate
it. A three-way ablation compares servoing with no weighting, with a purely
residual-based Tukey weighting, and with the proposed Hermite-informed
weighting, all else held fixed. The first variant is the scheme of
Chapter~\ref{ch:herpvs} exactly, so that the comparison of the first with the
second measures what robust estimation contributes, and the comparison of the
second with the third measures what the structural information contributes
over and above it.

%%% (5) Semicolons throughout, and "the results" now points at Section 4.7
%%% as a whole rather than at one of its three subsections.
The chapter is organised as follows. Section~\ref{sec:robust-framework}
introduces robust estimation and the weighted control law;
Section~\ref{sec:tukey} presents the Tukey weighting;
Section~\ref{sec:hermite-weight} develops the Hermite-informed weighting;
Section~\ref{sec:robust-stability} discusses stability;
Section~\ref{sec:ablation-protocol} defines the ablation protocol;
Section~\ref{sec:results} reports the results; and
Section~\ref{sec:ch4-conclusion} concludes.
%%% sec:results should label Section 4.7 itself. If that section carries no
%%% label at present, add \label{sec:results} to its heading.

% ======================================================================
%  OPTIONAL. If you would rather foreshadow the finding, which many
%  examiners prefer, append this to the ablation paragraph. It is accurate
%  as of Section 4.7.2 and commits to nothing that the runs do not show.
%  Do NOT use it unless the occluded results stand as they are.
%
%    It will be seen that the first comparison is decisive and the second is
%    not: robust weighting recovers the pose to within a quarter of a pixel
%    where the unweighted scheme does not recover it at all, while the two
%    weightings prove indistinguishable on the single compact occluder
%    tested here.
%
%  WHAT I DID NOT CHANGE
%  - the cite key26 for M-estimation, which I cannot verify from the PDF;
%  - \ref{sec:hermite-properties}, which resolved correctly to 3.2.4;
%  - the passive-tolerance paragraph, which is the strongest in the section
%    and is now supported by measurement rather than assertion.
% ======================================================================
